
\subsubsection{6.1 Interpretation of Asymmetry}

The observed asymmetry—distance diversity without structural diversity—can be attributed to the architectural design. The deterministic LSTM updates with a single learned initialization point enable the model to produce solutions at varying distances from ground truth (indicating differentiation of instance difficulty) while maintaining similar violation distribution patterns across instances.

\textbf{Distance Diversity Mechanism:} Message-passing aggregates constraint information through 32 rounds of LSTM updates, creating instance-specific embeddings. Different instances yield assignments at different distances from ground truth, demonstrating that the architecture differentiates problem characteristics.

\textbf{Structural Homogeneity Mechanism:} All instances begin from the same learned initialization and follow deterministic update rules. The architecture lacks mechanisms that would promote structural diversity:
\begin{itemize}
\item No stochastic sampling in message aggregation
\item No ensemble of initialization points
\item No diversity regularization in the loss function
\item No temperature-based exploration during decoding
\end{itemize}

The unsupervised classification loss optimizes for single-bit prediction accuracy rather than assignment quality or violation diversity. Training converged to the theoretical minimum, indicating that violation homogeneity persists under full optimization of this objective.

\subsubsection{6.2 Implications}

The finding suggests that baseline NeuroSAT does not generate the structural diversity hypothesized to support basin recovery stratification. Two-stage approaches that rely on heterogeneous violation patterns for effective refinement would require architectural modifications to Stage 1 before Stage 2 mechanisms could be tested.

Potential modifications include:
\begin{enumerate}
\item Stochastic message-passing with temperature-based sampling during aggregation
\item Ensemble initialization using multiple models with different random seeds
\item Diversity regularization with entropy maximization terms in the loss function
\end{enumerate}

These modifications directly address the deterministic convergence identified as the source of structural homogeneity.

\subsubsection{6.3 Limitations}

\textbf{Sample Size:} Only 8 satisfiable instances were evaluated after filtering 2 unsatisfiable instances from the original 10-instance test set. This limits statistical power for distribution analysis. Quartile-based range estimates have limited reliability with n=8; robust distribution statistics typically require n≥100.

\textbf{Training Scale:} Training used 20 instances rather than the full 80,000-instance G4SATBench dataset. While training converged to theoretical minimum loss, full-scale training might affect absolute metric values. However, the architectural characteristics (deterministic updates, single initialization) remain unchanged regardless of training scale.

\textbf{Single Seed:} One training run (seed=123) was conducted. Different random initializations might produce different entropy ranges. However, for identifying architectural limitations, a single proof-of-concept run provides sufficient evidence when the failure is substantial (entropy range 1.145 versus threshold 2.0).

\textbf{Difficulty Level:} Results are specific to 3-SAT easy instances (10-40 variables). Generalization to medium (40-200 variables) or hard (200-400 variables) difficulty levels has not been tested. The architectural mechanism (deterministic convergence to uniform strategies) is domain-independent, but empirical validation across difficulties remains future work.

\subsubsection{6.4 Methodological Contribution}

The dual-metric framework enables separate measurement of distance diversity and structural diversity, providing diagnostic information not available from single-metric evaluations (e.g., overall satisfaction rate). This approach may extend to other constraint satisfaction domains where characterizing solution diversity is relevant for multi-stage hybrid architectures.
